\begin{abstract}
    Este informe corresponde a la primera fase del Proyecto Integrador de Ciencia
    de Datos y plantea el problema de predecir el puntaje global de la prueba
    Saber 11 a partir de variables socioeconómicas del hogar, del colegio y del
    estudiante, formulado como una tarea de regresión. Se utiliza el conjunto
    público \textit{Resultados únicos Saber 11} del ICFES, restringido al periodo
    2022-4, descargado desde el Portal Nacional de Datos Abiertos. El análisis
    exploratorio revela que la mitad de las 1.065.888 filas descargadas son
    duplicados exactos, de modo que el conjunto de trabajo queda en 532.792
    estudiantes, y verifica que el puntaje global es una función determinista de
    los puntajes por área, que por tanto se excluyen como variables explicativas.
    La variable objetivo tiene media de 250,1 y desviación de 51,8 puntos, con
    una distribución casi simétrica. Las asociaciones más fuertes corresponden a
    la educación de la madre ($\rho = 0{,}35$), la edad ($\rho = -0{,}32$), la
    educación del padre (0,32) y la tenencia de computador (0,30) e internet
    (0,28); la jornada del colegio separa 70 puntos entre la completa y la
    nocturna, los colegios no oficiales superan a los oficiales en 29 puntos y la
    brecha entre departamentos alcanza 64 puntos entre Bogotá y Chocó. Una
    regresión lineal exploratoria explica el 24,2\,\% de la varianza con un error
    absoluto medio de 36,5 puntos, cota que servirá de referencia para las fases
    de modelado. El informe concluye con doce decisiones de preparación de datos
    e ingeniería de variables derivadas de los hallazgos.
\end{abstract}

\vspace{10pt}

\begin{IEEEkeywords}
    Ciencia de datos, análisis exploratorio de datos, Saber 11, ICFES, regresión,
    datos abiertos, fuga de información, nivel socioeconómico.
\end{IEEEkeywords}
